\chapter{Initial Data, Conservation, and the Action}

The laws of mechanics can be stated in two remarkably different ways. One
description is local: specify the state now, use a differential equation, and
advance the motion from one instant to the next. The other description treats
an entire trajectory as a single object and selects the physical path by a
stationary principle. Our purpose is to understand why both descriptions carry
the same physics.

We shall approach the global principle slowly. First we return to initial data
and the order of a differential equation. We then review energy and momentum,
not because their elementary Newtonian derivations are the deepest proofs of
those laws, but because they reveal structures that survive far beyond the
derivations. Only after that preparation shall we pass from ordinary functions
to functionals and uncover the action.

\section{What the Order of a Law Demands}

Begin in one dimension with a deliberately false law,
\begin{equation}
 F=m\dot x=m\frac{dx}{dt}.
 \label{eq:l02-aristotle}
\end{equation}
It says that a force is needed merely to sustain motion: remove the force and
the body stops at once. This is not Newtonian physics. Its value is precisely
that it lets us see what a first-order equation would demand of its initial
data.

Suppose the force is prescribed by position,
\begin{equation}
 F(x)=m\dot x.
 \label{eq:l02-first-order-force}
\end{equation}
If the position is known at one instant, then \(F(x)\) is known, and the law
immediately supplies the velocity,
\begin{equation}
 \dot x=\frac{F(x)}{m}.
\end{equation}
The same datum fixes the acceleration. Differentiating the force along the
motion gives
\begin{equation}
 \frac{dF}{dt}=\frac{dF}{dx}\dot x,
\end{equation}
and differentiating Eq.~\eqref{eq:l02-first-order-force} gives
\begin{equation}
 m\ddot x=\frac{dF}{dx}\dot x.
 \label{eq:l02-first-acceleration}
\end{equation}
Differentiate once more and use the product and chain rules:
\begin{align}
 m\dddot x
 &=\frac{d}{dt}\left(F'(x)\dot x\right) \\
 &=F''(x)\dot x^{\,2}+F'(x)\ddot x.
 \label{eq:l02-jerk}
\end{align}
Thus position determines velocity, velocity and position determine
acceleration, and the same chain continues through jerk and every higher time
derivative. In this first-order world, one number \(x(t_0)\) is enough to start
the motion.

Throughout this calculation the mass is constant. A visible subsystem may
seem to lose mass if a piece falls away, but then the discarded piece has
simply been omitted from the description. For the complete closed system, the
mass carried by the calculation has not disappeared.

Now replace the false law by Newton's law,
\begin{equation}
 F(x)=m\ddot x.
 \label{eq:l02-newton-one-d}
\end{equation}
Position still determines force and therefore acceleration, but it does not
determine velocity. Two bodies can pass through the same point with different
velocities and subsequently follow different motions. We must supply
\begin{equation}
 x(t_0)=x_0,
 \qquad
 \dot x(t_0)=v_0.
 \label{eq:l02-initial-data}
\end{equation}
Once these two data are known, Newton's equation supplies \(\ddot x\), and
successive differentiations supply the rest. The need for position and
velocity is not an arbitrary convention: it is the signature of a
second-order law.

For many particles, the statement simply expands. We must specify every
particle coordinate, every corresponding velocity component, and the force
law, which may depend on the relative positions of all the particles.

\begin{classroomqa}
\audiencequestion{Why does Newton's law require both position and velocity,
whereas the false first-order law requires only position?
\lecturetimestamp{00:09:14--00:11:57}}

\lecturerresponse{The distinction is the order of the equation. In
\(F(x)=m\dot x\), the position determines the first derivative directly. In
\(F(x)=m\ddot x\), it determines only the second derivative. The first
derivative remains independent and must be included in the initial state.}
\end{classroomqa}

\section{Energy and Conservative Forces}

Energy conservation is more general than the elementary argument we are about
to give. It remains meaningful in electrodynamics, special and general
relativity, and quantum mechanics. Here we ask the narrower question: under
what conditions does ordinary Newtonian particle mechanics conserve
\(T+U\)?

The qualification matters. A friction force appears to destroy mechanical
energy if the resulting heat and the microscopic degrees of freedom of the
environment are ignored. An even sharper counterexample is a tangential force
that continually pushes a constrained particle around a circle. After one
revolution the particle returns to the same position with a larger speed. Any
single-valued potential has returned to its original value, while the kinetic
energy has increased. Such a force cannot be written as the gradient of an
ordinary potential.

For the class of forces considered here, there is a potential energy
\(U(\mathbf x)\) such that
\begin{equation}
 F_i=-\frac{\partial U}{\partial x_i},
 \qquad
 \mathbf F=-\boldsymbol\nabla U.
 \label{eq:l02-conservative-force}
\end{equation}
The minus sign says that the force points toward decreasing potential energy.

Define
\begin{equation}
 T=\frac12 m\mathbf v^2
   =\frac12 m\sum_i v_i^2,
 \qquad
 E=T+U.
 \label{eq:l02-energy-definitions}
\end{equation}
The board develops the proof in the useful order: write the two energies,
differentiate them separately, and only then combine them.

\begin{figure}[H]
\centering
\includegraphics[width=0.68\textwidth]{lecture_02_figure_02.png}
\caption{The energy calculation at 00:22:24. The upper handwritten label is
read as \(T\): the formula, transcript, and later identity \(E=T+U\) fix that
interpretation.}
\lectureframe{00:22:24}
\label{fig:l02-energy-pieces}
\end{figure}

For the kinetic energy,
\begin{align}
 \frac{dT}{dt}
 &=\frac{d}{dt}\left(\frac12m\sum_i v_i^2\right) \\
 &=\sum_i m v_i\frac{dv_i}{dt}
 =\sum_i m v_i a_i.
 \label{eq:l02-tdot}
\end{align}
The potential changes because the particle moves through space. The
multivariable chain rule gives
\begin{align}
 \frac{dU}{dt}
 &=\sum_i\frac{\partial U}{\partial x_i}
             \frac{dx_i}{dt} \\
 &=\sum_i\frac{\partial U}{\partial x_i}v_i.
 \label{eq:l02-udot}
\end{align}
Adding the two derivatives,
\begin{align}
 \frac{dE}{dt}
 &=\sum_i v_i\left(ma_i+\frac{\partial U}{\partial x_i}\right) \\
 &=\sum_i v_i(F_i-F_i)=0.
 \label{eq:l02-energy-conservation}
\end{align}
The first \(F_i\) is Newton's \(ma_i\); the second comes from
\(\partial_iU=-F_i\). The cancellation is exact.

\begin{figure}[H]
\centering
\includegraphics[width=0.82\textwidth]{lecture_02_figure_03.png}
\caption{The completed energy derivative at 00:24:01. The right-hand bracket
is the decisive \(F_i-F_i\) cancellation.}
\lectureframe{00:24:01}
\label{fig:l02-energy-cancellation}
\end{figure}

Nothing in the argument restricts \(i\) to the three directions of one
particle. It may range over all coordinates of all particles. If each
generalized force is the negative derivative of one common potential and the
potential has no explicit time dependence, the total energy of the complete
system is conserved.

\begin{classroomqa}
\audiencequestion{Why denote potential energy by \(U\) rather than the common
letter \(V\)?\lecturetimestamp{00:18:13--00:18:36}}

\lecturerresponse{Either notation is possible. Here \(U\) avoids collisions:
\(V\) is often used for electric voltage and \(v\) already denotes velocity.
The symbol changes nothing in the physics.}
\end{classroomqa}

\section{Momentum and a Closed System}

Momentum enters by rewriting the same Newtonian equation. With constant mass,
\begin{equation}
 F_i=m\frac{dv_i}{dt}
     =\frac{d}{dt}(mv_i).
\end{equation}
Define
\begin{equation}
 \mathbf p=m\mathbf v.
\end{equation}
Then Newton's law becomes
\begin{equation}
 \mathbf F=\frac{d\mathbf p}{dt}.
 \label{eq:l02-force-momentum}
\end{equation}
In the absence of force, \(\dot{\mathbf p}=0\), so a free particle carries a
constant momentum. Intergalactic space is not perfectly empty, but it gives a
good approximation to this ideal.

\begin{figure}[H]
\centering
\includegraphics[width=0.62\textwidth]{lecture_02_figure_05.png}
\caption{The component, scalar, and vector forms of Newton's momentum law at
00:29:00.}
\lectureframe{00:29:00}
\label{fig:l02-momentum-law}
\end{figure}

This discussion uses a specifically Newtonian picture of time. Clocks can be
synchronized everywhere, all observers agree on one universal time, and time
flows uniformly. We are also considering force laws with no prescribed
external time dependence. These assumptions will not survive unchanged in
relativity, so they should be stated rather than hidden.

Now consider three interacting particles. Let \(\mathbf F_{ab}\) mean the
force on particle \(a\) due to particle \(b\). Internal forces give
\begin{align}
 \dot{\mathbf p}_1&=\mathbf F_{12}+\mathbf F_{13}, \\
 \dot{\mathbf p}_2&=\mathbf F_{21}+\mathbf F_{23}, \\
 \dot{\mathbf p}_3&=\mathbf F_{31}+\mathbf F_{32}.
 \label{eq:l02-three-particle-forces}
\end{align}
Newton's action--reaction law states
\begin{equation}
 \mathbf F_{ab}=-\mathbf F_{ba}.
 \label{eq:l02-third-law}
\end{equation}
The equal-and-opposite forces lie along the line joining the pair in the
elementary central-force picture, although only their cancellation is needed
here.

\begin{figure}[H]
\centering
\includegraphics[width=0.78\textwidth]{lecture_02_figure_04.png}
\caption{The board at 00:32:58, after the three force balances have been
assembled. The frame is used for their layout; the exact subscripts are
typeset in Eq.~\eqref{eq:l02-three-particle-forces}.}
\lectureframe{00:32:58}
\label{fig:l02-particle-balances}
\end{figure}

Adding the three equations gives
\begin{align}
 \frac{d}{dt}(\mathbf p_1+\mathbf p_2+\mathbf p_3)
 &=(\mathbf F_{12}+\mathbf F_{21})
  +(\mathbf F_{13}+\mathbf F_{31})\notag\\
 &\quad+(\mathbf F_{23}+\mathbf F_{32})\notag\\
 &=0.
\end{align}
Therefore
\begin{equation}
 \frac{d\mathbf P}{dt}=0,
 \qquad
 \mathbf P=\sum_a\mathbf p_a.
 \label{eq:l02-total-momentum}
\end{equation}

\begin{figure}[H]
\centering
\includegraphics[width=0.82\textwidth]{lecture_02_figure_06.png}
\caption{The completed cancellation board at 00:35:20. Each internal force
appears once with each orientation, leaving \(\dot{\mathbf P}=0\).}
\lectureframe{00:35:20}
\label{fig:l02-momentum-cancellation-board}
\end{figure}

\begin{figure}[htbp]
\centering
\begin{tikzpicture}[>=Stealth,thick,scale=1.05]
  \node[circle,fill=black,inner sep=1.8pt,label=left:$1$] (p1) at (-2.5,0) {};
  \node[circle,fill=black,inner sep=1.8pt,label=above:$2$] (p2) at (1.8,1.65) {};
  \node[circle,fill=black,inner sep=1.8pt,label=below:$3$] (p3) at (1.8,-1.65) {};
  \draw[->,blue] ($(p1)!0.08!(p2)$) -- ($(p1)!0.43!(p2)$)
    node[midway,above] {$\mathbf F_{12}$};
  \draw[->,blue] ($(p2)!0.08!(p1)$) -- ($(p2)!0.43!(p1)$)
    node[midway,below] {$\mathbf F_{21}$};
  \draw[->,red] ($(p1)!0.08!(p3)$) -- ($(p1)!0.43!(p3)$)
    node[midway,below] {$\mathbf F_{13}$};
  \draw[->,red] ($(p3)!0.08!(p1)$) -- ($(p3)!0.43!(p1)$)
    node[midway,above] {$\mathbf F_{31}$};
  \draw[->,teal] ($(p2)!0.12!(p3)$) -- ($(p2)!0.42!(p3)$)
    node[midway,right] {$\mathbf F_{23}$};
  \draw[->,teal] ($(p3)!0.12!(p2)$) -- ($(p3)!0.42!(p2)$)
    node[midway,left] {$\mathbf F_{32}$};
\end{tikzpicture}
\caption{Pairwise internal forces cancel when the momentum equations for a
closed system are summed.}
\label{fig:l02-pairwise-redraw}
\end{figure}

A closed system is one for which every relevant force source has been included.
If an external force is present, it does not acquire a partner inside the sum,
and the result becomes \(\dot{\mathbf P}=\mathbf F_{\mathrm{ext}}\).

\begin{classroomqa}
\audiencequestion{Are action and reaction instantaneous?
\lecturetimestamp{00:36:15--00:37:12}}

\lecturerresponse{They are instantaneous in Newtonian mechanics, where the
speed at which influences propagate is effectively infinite. Real influences
do not propagate instantaneously; special relativity forbids that. Yet
momentum conservation survives. The Newtonian derivation is therefore less
general than the conservation law it reveals.}
\end{classroomqa}

Momentum conservation applies to radiation pressure, quantum systems, and
relativistic physics, even though the simple pair-force proof cannot be carried
unchanged into every one of those settings. There are three momentum laws, one
for each spatial component, and one energy law. In special relativity these
four quantities belong together as energy--momentum. This is a first hint that
the conserved quantities have a deeper origin than the present calculations.

\begin{classroomqa}
\audiencequestion{Why does momentum seem to provide more conservation laws
than energy?\lecturetimestamp{00:37:37--00:38:24}}

\lecturerresponse{Momentum has one conserved component for each independent
direction of space, whereas energy is one scalar associated with evolution in
time. In three spatial dimensions that gives three momentum components and one
energy. Relativity later unifies the four rather than treating them as an
accidental list.}
\end{classroomqa}

\section{Stationarity in Ordinary Calculus}

We now prepare for what is really the central law of the course: the principle
of least action. The required mathematics begins with an ordinary function of
one variable, \(f(y)\). At a smooth local minimum,
\begin{equation}
 \frac{df}{dy}=0.
 \label{eq:l02-one-variable-stationary}
\end{equation}
The derivative also vanishes at a smooth maximum. It can even vanish at a flat
inflection point, such as the origin of \(f(y)=y^3\). The first-derivative
condition therefore identifies a \emph{stationary point}; the second derivative
and higher information distinguish the possibilities.

\begin{figure}[H]
\centering
\includegraphics[width=0.78\textwidth]{lecture_02_figure_07.png}
\caption{Minimum, maximum, and flat inflection behavior on the board at
00:43:20. All are stationary to first order.}
\lectureframe{00:43:20}
\label{fig:l02-stationary-board}
\end{figure}

For a function of many variables, stationarity must hold along every coordinate
direction:
\begin{equation}
 \frac{\partial f}{\partial x_i}=0
 \qquad\text{for every }i.
 \label{eq:l02-many-stationary}
\end{equation}
Geometrically, at the bottom of a bowl there is no downhill first-order
direction. The same equations also find maxima and saddles, so ``stationary''
remains the accurate word.

As an example, take
\begin{equation}
 f(x,y)=x^2+y^2+xy+3x+2y+7.
 \label{eq:l02-quadratic-example}
\end{equation}
The stationarity equations are
\begin{align}
 \frac{\partial f}{\partial x}&=2x+y+3=0, \\
 \frac{\partial f}{\partial y}&=x+2y+2=0.
 \label{eq:l02-quadratic-stationarity}
\end{align}
Solving the linear system gives
\begin{equation}
 x=-\frac43,
 \qquad
 y=-\frac13.
 \label{eq:l02-correct-stationary-point}
\end{equation}
The constant \(7\) never enters the equations because a vertical shift changes
the value of \(f\), not the location of its stationary point.

\begin{figure}[H]
\centering
\includegraphics[width=0.84\textwidth]{lecture_02_figure_08.png}
\caption{The two-variable example at 00:51:10. The displayed derivative
equations are correct; the board's \(x=-8/3\) is an arithmetic slip. Solving
those equations gives Eq.~\eqref{eq:l02-correct-stationary-point}.}
\lectureframe{00:51:10}
\label{fig:l02-quadratic-board}
\end{figure}

For \(n\) variables there are \(n\) first-order equations. They may have no
solution, one solution, or several solutions. A monotone function need not
have a stationary point, while a landscape with several valleys can have many.

\begin{classroomqa}
\audiencequestion{Is the physical principle really one of least action, or of
stationary action?\lecturetimestamp{00:42:05--00:46:39}}

\lecturerresponse{Stationary action is the precise statement. The first-order
change of the action vanishes; that does not by itself prove a minimum. In many
elementary mechanical problems the physical path is locally minimizing, so
``least action'' is useful language, but the distinction should not be
forgotten. Velocity-dependent interactions are one reason not to build the
theory around a naive minimum test.}
\end{classroomqa}

\section{Generalized Trajectories}

The basic mathematical problem of mechanics is to determine the trajectory of
a system from suitable data. For \(10^{23}\) particles, no one attempts to
write every trajectory in practice; statistical or approximate descriptions
are used instead. Nevertheless, the exact classical question remains well
defined.

Describe a system by generalized coordinates
\begin{equation}
 q_1,q_2,\ldots,q_N.
\end{equation}
They may be Cartesian positions, angles, orientations, latitude and longitude,
or any other independent variables that locate the system in configuration
space. A single particle in three-dimensional space has three such
coordinates; two particles have six; \(n\) unconstrained particles have
\(3n\).

The motion is the whole collection
\begin{equation}
 \mathbf q(t)=\bigl(q_1(t),q_2(t),\ldots,q_N(t)\bigr).
 \label{eq:l02-generalized-trajectory}
\end{equation}
This curve through configuration space is a generalized trajectory. To begin a
Newtonian prediction, we specify both \(\mathbf q(t_0)\) and
\(\dot{\mathbf q}(t_0)\). Newton's equations then act locally: knowing where the
system is and how it is moving tells us its acceleration, and hence how to
advance to the neighboring point of the path.

\begin{figure}[H]
\centering
\includegraphics[width=0.82\textwidth]{lecture_02_figure_09.png}
\caption{The generalized trajectory at 01:00:20. The horizontal plane stands
for the many-dimensional configuration coordinates; the vertical direction
represents time.}
\lectureframe{01:00:20}
\label{fig:l02-generalized-trajectory-board}
\end{figure}

There is another way to pose the problem. Instead of giving position and
velocity at one time, hold the configurations fixed at two times,
\begin{equation}
 \mathbf q(t_0)=\mathbf q_0,
 \qquad
 \mathbf q(t_1)=\mathbf q_1,
 \label{eq:l02-fixed-endpoints}
\end{equation}
and ask which trajectory connects them. In the elementary well-posed cases
being considered, this supplies the same amount of boundary data as position
and velocity at one time. More complicated systems or long time intervals can
admit several classical paths, but each stationary path still satisfies the
same equations of motion.

The contrast is conceptual. Newton's differential equation gives a local rule
along the curve. A variational principle assigns one number to the whole curve
and asks which admissible curve makes that number stationary.

\section{From Functions to Functionals}

Take two fixed points in the \(xy\)-plane and describe a connecting curve by
\(y(x)\). The shortest connection is plainly a straight line, but we want the
mathematical question that produces that answer.

For a small segment,
\begin{equation}
 ds^2=dx^2+dy^2,
\end{equation}
and therefore
\begin{equation}
 ds=dx\sqrt{1+\left(\frac{dy}{dx}\right)^2}.
 \label{eq:l02-line-element}
\end{equation}
Adding all the little segments gives the length
\begin{equation}
 \mathcal S[y]
 =\int_{x_1}^{x_2}
   \sqrt{1+\left(y'(x)\right)^2}\,dx.
 \label{eq:l02-arc-length-functional}
\end{equation}

\begin{figure}[H]
\centering
\includegraphics[width=0.82\textwidth]{lecture_02_figure_10.png}
\caption{The curve, its \(dx\)-\(dy\) triangle, and the arc-length integral at
01:12:00.}
\lectureframe{01:12:00}
\label{fig:l02-arc-length-board}
\end{figure}

The input to \(\mathcal S\) is not one number or a finite list of numbers. It is
the entire function \(y(x)\), and the output is a number. Such an object is a
\emph{functional}.

There are complementary global and local views. Globally, compare every curve
between the endpoints and select a shortest one. Locally, break the curve into
tiny pieces and focus on a point between two nearby pinned neighbors. Moving
that point until the two small segments line up says simply: continue straight
ahead. Applying the same local condition everywhere gives the straight line.
The global extremization and the local differential equation encode the same
geometry.

One tempting shortcut fails. Since the integrand in
Eq.~\eqref{eq:l02-arc-length-functional} is smallest when \(y'=0\), why not set
\(y'=0\) at every point? Because a horizontal line will not generally meet the
fixed endpoints. The admissible curves are constrained globally. Minimizing an
integrand point by point while ignoring its boundary conditions solves a
different problem.

\subsection{Fermat's Principle}

Optics supplies a more revealing functional. Suppose a layered isotropic
medium has a light speed \(c(y)\) that varies with height but not with direction
at a fixed point. Fermat's principle says that the physical ray makes the
travel time stationary. Locally,
\begin{equation}
 c(y)=\frac{ds}{dt},
 \qquad
 dt=\frac{ds}{c(y)}.
\end{equation}
Using Eq.~\eqref{eq:l02-line-element},
\begin{equation}
 dt=\frac{\sqrt{1+(y')^2}}{c(y)}\,dx,
\end{equation}
so the total time is the functional
\begin{equation}
 \mathcal T[y]
 =\int_{x_1}^{x_2}
 \frac{\sqrt{1+\left(y'(x)\right)^2}}{c(y)}\,dx.
 \label{eq:l02-fermat-functional}
\end{equation}

\begin{figure}[H]
\centering
\includegraphics[width=0.84\textwidth]{lecture_02_figure_11.png}
\caption{The layered-medium ray and completed least-time functional at
01:23:10. The ray bends toward the faster region while still advancing toward
its endpoint.}
\lectureframe{01:23:10}
\label{fig:l02-fermat-board}
\end{figure}

If \(c\) is constant, minimizing time is the same as minimizing distance, and
the ray is straight. If light is faster higher in the medium, the best path can
curve upward to spend more of its journey in the faster region. It cannot gain
indefinitely by going straight up and then straight across: the vertical detour
costs time without advancing toward the destination. The physical path is a
compromise selected by the whole integral.

\subsection{The Hanging Chain}

A chain of fixed length suspended from two points supplies another variational
problem. Gravity drives each small mass element downward, while the fixed
length prevents every element from falling independently. The equilibrium
curve minimizes the total gravitational potential energy subject to the length
constraint.

\begin{figure}[H]
\centering
\includegraphics[width=0.68\textwidth]{lecture_02_figure_12.png}
\caption{The hanging-chain sketch at 01:25:20. The fixed endpoints and fixed
length constrain the chain while gravity lowers its center of mass.}
\lectureframe{01:25:20}
\label{fig:l02-catenary-board}
\end{figure}

The resulting curve is a catenary, which may be written
\begin{equation}
 y(x)=a\cosh\!\left(\frac{x-x_0}{a}\right)+y_0
\end{equation}
after the constants are chosen to fit the endpoints and length. Its derivation
is not needed yet; its role is to show that minimizing a functional is a common
physical task, not an artificial game invented for mechanics.

\begin{classroomqa}
\audiencequestion{From what height is the chain's potential energy measured?
\lecturetimestamp{01:25:50--01:26:16}}

\lecturerresponse{Any fixed reference height will do. Shifting the zero adds
the same constant to the energy of every allowed configuration, so it cannot
change which configuration minimizes the energy.}
\end{classroomqa}

These examples define the calculus of variations: find a function that makes a
functional stationary. Mechanics will ask the same question for all the
functions \(q_i(t)\) that make up a trajectory.

\section{The Action and the Lagrangian}

We can now state the global formulation of mechanics. For every mechanical
system there is a quantity called the action. Hold the endpoint configurations
fixed and compare neighboring trajectories between them. The physical
trajectory makes the action stationary.

\begin{classroomqa}
\audiencequestion{Is this really equivalent to Newton's formulation, or merely
an approximation to it?\lecturetimestamp{01:28:21--01:30:13}}

\lecturerresponse{For systems governed by Newtonian mechanics, it is an
equivalent formulation: varying the correct action yields Newton's equations.
Which language is more efficient depends on the problem. The action principle
also extends naturally to relativistic particles, fields, waves, and
electromagnetism, where a literal catalogue of pairwise Newtonian forces is no
longer the right description.}
\end{classroomqa}

The practical advantage can be dramatic. Imagine describing a rigid eraser by
writing force equations for every microscopic constituent. The local equations
exist in principle, but they hide the small set of collective variables we
actually care about: translation and rotation. A suitable action written in
those variables reaches the same physics far more directly. Likewise, in a
frictionless block-on-wedge problem, conservation laws or a Lagrangian can
avoid pages of constraint forces. A more general principle is often not only
deeper but faster.

For one particle moving in one dimension, the unknown path is \(x(t)\). Since
kinetic and potential energy have dominated the discussion, a natural first
guess is
\begin{equation}
 \int_{t_0}^{t_1}(T+U)\,dt.
\end{equation}
That is the total energy integrated over time, and it is the wrong guess. The
correct action is
\begin{align}
 A[x]
 &=\int_{t_0}^{t_1}(T-U)\,dt \\
 &=\int_{t_0}^{t_1}
   \left[\frac12m\dot x^{\,2}-U(x)\right]dt.
 \label{eq:l02-action-one-particle}
\end{align}
The difference
\begin{equation}
 L(x,\dot x)=T-U
 \label{eq:l02-lagrangian-one-particle}
\end{equation}
is the Lagrangian, and the action is its integral over time.

The minus sign is not a decorative convention. If \(U\) is defined by
\(F=-dU/dx\), varying \(T-U\) produces the correct sign in Newton's equation.
One could rename \(-U\) as a new quantity and write a plus sign, but that new
quantity would then be the negative of potential energy. Renaming the symbol
does not change the relation.

\begin{figure}[H]
\centering
\includegraphics[width=0.84\textwidth]{lecture_02_figure_13.png}
\caption{The action and its general Lagrangian form at 01:43:20. The
independent variable is now time, and the path is the collection of functions
\(q_i(t)\).}
\lectureframe{01:43:20}
\label{fig:l02-action-board}
\end{figure}

For generalized coordinates, write
\begin{equation}
 L=L(\mathbf q,\dot{\mathbf q}),
 \qquad
 A[\mathbf q]
 =\int_{t_0}^{t_1}L(\mathbf q,\dot{\mathbf q})\,dt.
 \label{eq:l02-general-action}
\end{equation}
The analogy with the earlier examples is exact at the structural level. The
arc-length integrand depends on \(y\) and \(y'\); the optical integrand depends
on \(y\), \(y'\), and \(c(y)\); the mechanical integrand depends on
\(\mathbf q\) and \(\dot{\mathbf q}\). In every case an integral assigns a
number to a whole path with fixed endpoints.

For ordinary Newtonian particles, \(L=T-U\). More general classical systems
can have less familiar Lagrangians, especially when velocity-dependent
interactions or fields are present. What survives is the architecture:
construct the action, vary the path while holding its endpoints fixed, and set
the first-order change to zero.

\begin{classroomqa}
\audiencequestion{Why kinetic energy minus potential energy, rather than the
total energy?\lecturetimestamp{01:33:30--01:36:38}}

\lecturerresponse{At this stage the honest test is operational. Use the
calculus of variations and see which choice reproduces Newton's equations.
\(T-U\) does; \(T+U\) gives the wrong force sign. The derivation, rather than an
analogy with conserved energy, will explain why this is the right combination.}
\end{classroomqa}

That derivation is the next task. Given a Lagrangian, we must learn how to turn
stationary action into local equations of motion, and how to move back and
forth between those equations and Newton's law. The same technique will solve
the optical least-time problem. Mechanics and geometrical optics have arrived
at the same mathematical threshold from different directions.
